%% paginas/4_pos_textuais/glossario.tex
%% Elemento Pós-Textual Opcional (Glossário) -- Conforme ABNT NBR 14724
%% Regra: Relação alfabética de termos técnicos acompanhados de suas definições.

\chapter*{Glossário}
\addcontentsline{toc}{chapter}{GLOSSÁRIO} % Adiciona de forma limpa ao sumário

\begin{spacing}{1.5}
  \noindent
  \textbf{Compilação}: Processo de conversão do código fonte em linguagem de marcação LaTeX para um arquivo final visualizável, tipicamente em formato PDF.
  
  \vspace{0.5em}
  \noindent
  \textbf{Estilo modular}: Abordagem de arquitetura onde as regras de visualização e design são separadas logicamente da matéria escrita, organizando-se o projeto em múltiplos arquivos pequenos e interdependentes.
  
  \vspace{0.5em}
  \noindent
  \textbf{Overleaf}: Plataforma de edição colaborativa online de código LaTeX, amplamente utilizada no meio acadêmico e científico para coautoria de relatórios, artigos e teses.
  
  \vspace{0.5em}
  \noindent
  \textbf{Pseudocódigo}: Método de descrição de algoritmos que utiliza uma mistura de linguagem natural e estruturas de programação de forma simplificada, facilitando o entendimento lógico sem a rigidez sintática de uma linguagem de programação real.
\end{spacing}
